\section{低秩近似：用少数方向保留主要结构}
\label{sec:03-Linear-Algebra/07-SVD-and-Data-Applications/03}

\subsection{数据往往不是低秩的，却可能接近低秩}

SVD 把一个矩阵写成了秩一成分的和：

\[
A = \sum_{i=1}^{r} \sigma_i u_i v_i^{\trans}.
\]

每一项做一件简单的事：读取输入在 \(v_i\) 方向上的坐标，乘以 \(\sigma_i\)，沿 \(u_i\) 方向输出。第一项贡献最多能量，第二项次之，依此类推。

现在翻过这张图来看数据。你手上有一个 \(1000\times 1000\) 的矩阵——比如用户-物品评分表，或者图像像素阵列。一百万个数字。你对它做 SVD，发现前 20 个奇异值合计占了总能量的 95\%，从第 21 个开始，每个奇异值都小到像是噪声。数学上，这个矩阵的秩可能等于 1000——毕竟每一行都有独特的测量误差或像素值的微小波动，从线性无关性的严格定义来说，它们确实彼此独立。但它的"有效信息维数"大概只有 20。

这就引出一个问题：如果只保留前 \(k\) 个秩一成分，扔掉其余，误差有多大？更关键的是：在所有秩不超过 \(k\) 的矩阵中，这个截断版本是不是最好的？

\subsection{截断 SVD}

将奇异值从大到小排列：\(\sigma_1\ge\cdots\ge\sigma_r>0\)。取前 \(k\) 项：

\[
A_k = \sum_{i=1}^{k} \sigma_i u_i v_i^{\trans}
= U_k \Sigma_k V_k^{\trans}.
\]

\(A_k\) 的秩至多为 \(k\)。其余 \(r-k\) 个奇异方向被整体丢弃。差矩阵恰好是剩下的部分：

\[
A - A_k = \sum_{i=k+1}^{r} \sigma_i u_i v_i^{\trans}.
\]

因为被丢弃的部分本身也有 SVD 的结构（基向量正交、奇异值排序），两个常用的误差度量可以立刻读出：

\[
\|A - A_k\|_2 = \sigma_{k+1},
\qquad
\|A - A_k\|_F^2 = \sum_{i=k+1}^{r} \sigma_i^2.
\]

谱范数 \(\|\cdot\|_2\) 报告最坏输入方向上的误差——存在某个单位向量，经 \(A\) 和 \(A_k\) 作用后的差别正好是 \(\sigma_{k+1}\)。Frobenius 范数则把所有被舍弃方向的平方误差加总，报告"总体"偏差。

两者的关注点不同。谱范数适合最坏情形分析：你最关心的那个输入，近似之后偏了多少？Frobenius 范数适合总体质量评估：整个矩阵被削掉了多少能量？

\subsection{为什么它是最好的：Eckart--Young 定理}

Eckart--Young 定理给出了一个明确答案：在谱范数或 Frobenius 范数下，没有任何秩不超过 \(k\) 的矩阵能比 \(A_k\) 更接近 \(A\)。

谱范数情形的证明有一个漂亮的维数论据。任取秩不超过 \(k\) 的矩阵 \(B\)。由秩--零化度定理，\(B\) 的零空间维数至少为 \(n-k\)。另一方面，\(v_1,\ldots,v_{k+1}\) 张成的子空间维数为 \(k+1\)（这些是前 \(k+1\) 个右奇异向量）。两个子空间的维数之和是 \((n-k)+(k+1)=n+1>n\)，所以它们的交必非零：存在单位向量 \(x\)，同时落在 \(\operatorname{span}\{v_1,\ldots,v_{k+1}\}\) 和 \(N(B)\) 中。

\(Bx=0\)，于是

\[
\|(A-B)x\| = \|Ax\|.
\]

把 \(x\) 写为 \(x=\sum_{i=1}^{k+1}c_i v_i\)，其中 \(\sum c_i^2=1\)（因为是单位向量）。则

\[
\|Ax\|^2 = \Bigl\|\sum_{i=1}^{k+1} c_i Av_i\Bigr\|^2
= \Bigl\|\sum_{i=1}^{k+1} c_i\sigma_i u_i\Bigr\|^2
= \sum_{i=1}^{k+1} c_i^2\sigma_i^2.
\]

所有 \(\sigma_i\ge\sigma_{k+1}\)（奇异值递减），所以

\[
\|Ax\|^2 \ge \sigma_{k+1}^2 \sum_{i=1}^{k+1} c_i^2 = \sigma_{k+1}^2.
\]

因此存在单位向量 \(x\) 使得 \(\|(A-B)x\|\ge\sigma_{k+1}\)，即 \(\|A-B\|_2\ge\sigma_{k+1}\)。而 \(A_k\) 恰好达到这个下界（误差就是 \(\sigma_{k+1}\)），所以它是最优的。

Frobenius 范数的结论来自正交不变性：Frobenius 范数在正交变换下不变，因此任何秩 \(k\) 近似的误差平方至少是前 \(k\) 个方向"没被捕捉到"的平方能量。精确计算表明，所有秩 \(k\) 矩阵的 Frobenius 误差下界都是 \(\sum_{i>k}\sigma_i^2\)，\(A_k\) 再次恰好达到。

这个定理的力量在于它的前提极少：只要求秩不超过 \(k\)，不要求矩阵对称、稀疏或非负。但正因前提如此之少，它的结论也仅限于此。如果近似值还必须满足非负性（比如图像像素值）、稀疏性（比如网络邻接矩阵）、整数性（比如计数矩阵）或物理守恒律，截断 SVD 可能不再可行，更不一定是最优的。定理说的是"在只有秩约束时"的最优解，不是所有约束下的通用答案。

\subsection{一个可以算完的压缩例子}

设某矩阵的奇异值为 \((10, 3, 0.2)\)。取秩二近似后，最坏方向误差是 \(0.2\)，Frobenius 误差也是 \(0.2\)（只剩下一个方向）。保留的平方能量比例为

\[
\frac{10^2 + 3^2}{10^2 + 3^2 + 0.2^2}
= \frac{109}{109.04} \approx 99.963\%.
\]

存储开销的变化同样可以精确计算。原矩阵 \(m\times n\) 直接存需要 \(mn\) 个数。秩 \(k\) 截断 SVD 的因子形式 \(U_k\Sigma_k V_k^{\trans}\) 只需储存 \(U_k\)（\(m\times k\)）、\(\Sigma_k\) 的对角元（\(k\) 个）和 \(V_k\)（\(n\times k\)），合计约 \(k(m+n+1)\) 个数。当 \(k\) 远小于 \(m\) 和 \(n\) 时，这种表示同时节省了存储空间和矩阵-向量乘法的计算量。

但这里有一个容易踩的坑。"保留了 99.963\% 能量"听起来像是"保留了 99.963\% 信息"。奇异值只描述矩阵在欧氏几何中的长度-角度结构。一个数值极小的奇异方向，可能恰好携带分类边界所需的那一丁点差异信号；一个数值极大的奇异方向，也可能只是全局亮度或量纲差异带来的背景。奇异值告诉你一个方向在矩阵的作用下放了多大，不告诉你这个方向对你的任务有多大意义。前者是数学，后者是语义。

\subsection{从去噪到缺失数据}

如果后部奇异方向携带的主要是各向同性的小噪声，截断 SVD 会产生自然的去噪效果——噪声分散在大量小奇异值上，信号集中在少数大奇异值上，截断刚好把噪声剥离。这是 SVD 去噪的基本逻辑。

但这个逻辑依赖一个没有写明的假设：噪声和信号的奇异向量大致对齐。如果噪声恰好落在某个大奇异值的方向上，截断不仅不去噪，反而帮倒忙。奇异值本身不携带"噪声/信号"标签。选择 \(k\) 时，奇值衰减曲线上的"肘部"提供了一个视觉参考点，但不应是唯一依据。误差预算、下游任务指标和领域知识通常更可靠。

缺失数据又加了一层复杂度。普通 SVD 假定矩阵的每个元素都已知。如果矩阵有缺失项——比如用户只评过少数物品——把缺失项粗暴填成零，等于把"未知"当成了"确知为零"。这会对 SVD 的低秩结构产生系统性的偏误。

矩阵补全之所以可能，是因为它额外假设了矩阵潜在低秩，并且观测位置足够分散（不能只观测某几行或某几列）。低秩假设告诉你矩阵的列（或行）之间存在冗余——知道了足够多的元素，其余元素可以通过这种冗余推断出来。但这不是魔法。如果只观测了前 10 行，你对第 11 行一无所知——低秩假设无法从一片空白中猜测出你没有见过的行。低秩是信息恢复的结构性先验，不是包罗万象的万能补丁。

如果追求的是核范数最小化或矩阵补全的凸松弛解法，需要先理解本节的基本结构：截断 SVD 是在完整数据上、仅受秩约束时的最优近似。加上缺失、噪声结构或额外约束后，最优解的形态会改变，但截断 SVD 总是比较的基准。

\subsection{延伸与来源}

\begin{itemize}
\tightlist
\item
  MIT 18.065：Low Rank Matrices and Matrix Factorization
\item
  MIT 18.06SC：Singular Value Decomposition
\end{itemize}
